\documentclass{article}
\usepackage{amsmath} 
\usepackage{amssymb}
\usepackage[a4paper]{geometry}
\usepackage{fancyhdr}
\pagestyle{fancy}
\lhead{Der Einzelspalt}
\rhead{März 2026}
\begin{document}
\section{Der Einzelspalt} 
Der Einfachspalt, auch Einfachspalt genannt, besteht aus einem Spalt, durch den ein Laser auf ein Schirm gestrahlt wird.
 
\subsection{Das Interferenzmuster}
Das Interferenzmuster eines Einfachspalt kann begründet werden, indem das durchkommende Licht in eine Anzahl $n$ Zonen aufgeteilt wird.
 
Für einen jeden Winkel der das Licht haben kann, wenn es den Spalt verlässt, kann das Licht in eine Anzahl an $n$ Zonen aufgeteilt werden, so dass das Licht der einen Zone zur nächsten eine Phasendifferenz von $\lambda / 2$ hat. Ist nun $n = 2k$ mit $k \in \mathbb{N}$, so kann jeder Zone eine andere Zone zugeteilt werden, so dass alle Zonen gegenseitig destruktiv interferieren. Ist $n = 2k + 1$, so gibt es nach der destruktiven Interferenz eine weitere Zone, wessen Licht den Schirm treffen kann.
 
Der Anteil des gesamten Lichtes welcher den Schirm trifft ist also $1/n$. 
\end{document}